\documentclass[11pt]{amsart}
\usepackage[T1]{fontenc}
\usepackage{lmodern}
\usepackage{amsmath,amssymb,amsthm,mathtools}
\usepackage{booktabs,array,longtable}
\usepackage{microtype}
\usepackage{needspace}
\usepackage[hidelinks]{hyperref}
\hypersetup{pdftitle={The L'vov-Kaplansky Conjecture for Standard Polynomials},pdfauthor={Y. T. Miah},pdfsubject={Set-of-values classification for alternating standard polynomials}}
\usepackage{enumitem}
\newtheorem{theorem}{Theorem}[section]
\newtheorem{proposition}[theorem]{Proposition}
\newtheorem{lemma}[theorem]{Lemma}
\newtheorem{corollary}[theorem]{Corollary}
\newtheorem{remark}[theorem]{Remark}
\newcommand{\Mat}{\mathrm{M}}
\newcommand{\im}{\operatorname{Im}}
\newcommand{\tr}{\operatorname{tr}}
\newcommand{\diag}{\operatorname{diag}}
\newcommand{\Span}{\operatorname{span}}
\newcommand{\rk}{\operatorname{rank}}
\newcommand{\one}{\mathbf{1}}
\newcommand{\ord}{\operatorname{ord}}
\newcommand{\Sl}{\mathfrak{sl}}
\newcommand{\Gam}{\Gamma}
\title[L'vov--Kaplansky for Standard Polynomials]{The L'vov--Kaplansky Conjecture for Standard Polynomials}
\author{Y. T. Miah}
\date{7 November 2026}
\subjclass[2020]{16R10, 15A24, 16R30}
\keywords{L'vov--Kaplansky conjecture, standard polynomial, matrix polynomial image, alternating polynomial, Euler trail, generic division algebra}
\emergencystretch=2em
\begin{document}
\begin{abstract}
The L'vov--Kaplansky conjecture asks whether the image of a multilinear
polynomial on a full matrix algebra is a vector subspace. For the
unnormalized standard polynomial $s_d$ over an algebraically closed field
$K$ of characteristic zero, the sets of individual matrix values are
determined in every degree and order. The image of $s_{2r-1}$ is zero
when $n<r$ and all of $\Mat_n(K)$ when $n\ge r$; the image of $s_{2r}$
is zero when $n\le r$ and $\Sl_n(K)$ when $n\ge r+1$.
The odd-degree proof uses explicit linear slices and Jordan forms.
The even-degree argument combines a uniform nilpotent construction and
an additive transfer for admissible cores with generic-nullity results
of Brassil--Reichstein and Lin--Tang--Zhao, extending a method of
Qiyue Tang to obtain actual polynomial values.
\end{abstract}
\maketitle

\section{Introduction and the classification}\label{sec:introduction}
The L'vov--Kaplansky conjecture concerns the images of multilinear
noncommutative polynomials evaluated on full matrix algebras. In the
setting of algebraically closed fields of characteristic zero, it asks
whether each such image is a vector subspace
\cite{KanelBelovMalevRowen2016,KanelBelovEtAl2020}.
The alternating standard polynomials constitute a distinguished family
for which one can go further and identify the image explicitly, in all
degrees and matrix orders.

For every integer $d\ge1$ define the \emph{unnormalized standard polynomial}
\begin{equation}\label{eq:standard}
 s_d(X_1,\ldots,X_d)=\sum_{\sigma\in S_d}\operatorname{sgn}(\sigma)
 X_{\sigma(1)}\cdots X_{\sigma(d)}.
\end{equation}
The image $\im s_d(\Mat_n(K))$ means the set of matrices attained as
$s_d(A_1,\ldots,A_d)$ for \emph{one tuple} of inputs; it is not the
linear span or Zariski closure of those values. Write
$\Sl_n(K)=\{T\in\Mat_n(K):\tr T=0\}$. Matrix units are indexed from
zero except when quoting another author's conventions. Throughout the
article, $K$ is algebraically closed of characteristic zero unless a
statement explicitly declares otherwise.

\Needspace{13\baselineskip}
\begin{theorem}[Images of the standard polynomials]\label{thm:main}
For every algebraically closed field $K$ of characteristic zero and all
integers $r,n\ge1$, the \emph{actual images} are
\begin{equation}\label{eq:main}
\begin{aligned}
 \im s_{2r-1}(\Mat_n(K))
 &=\begin{cases}\{0\},&n<r,\\ \Mat_n(K),&n\ge r,\end{cases}\\[3pt]
 \im s_{2r}(\Mat_n(K))
 &=\begin{cases}\{0\},&n\le r,\\ \Sl_n(K),&n\ge r+1.\end{cases}
\end{aligned}
\end{equation}
\end{theorem}

The zero branches follow from the Amitsur--Levitzki identity and
its standard recursion \cite{AmitsurLevitzki1950}. In even degree,
cyclically rotating a word reverses its permutation sign while
preserving its trace, so every value is traceless. The main task is
therefore to attain each matrix in the two positive branches of
\eqref{eq:main}.

The lowest-degree cases are well established. The ordinary commutator
$s_2$ attains every traceless matrix \cite{AlbertMuckenhoupt1957},
and Vitas's cubic image-linearity theorem gives the corresponding
$s_3$ result \cite{Vitas2026}. The quartic result of Fagundes,
de Mello and da Silva dos Santos includes the actual-image assertion
for $s_4$ \cite[Theorem~8, Case~1]{FagundesMelloSantos2022}.
Qiyue Tang established the degree-five and degree-six conclusions,
including the actual traceless image of $s_6$ in orders at least four
\cite[Appendix~A, Proposition~11]{Tang2026FiveSix}.

Classical antecedents include Kostant's work on standard identities,
Swan's graph method and its correction, and results of Rowen and
Hutchinson on standard polynomials and matrix identities
\cite{Kostant1958,Swan1963,Swan1969Correction,Rowen1974,Hutchinson1974}.
Serre studied special matrix values and related norm questions
\cite{Serre2016}. These works provide context for the exact
all-degree image classification undertaken here.

The two parity arguments have different structures. Part~I establishes
the odd-degree assertion through signed path and cycle evaluations,
fixed linear slices and target-dependent Jordan constructions at the
minimal size. Part~II uses a rank-one nilpotent with surjective full
differential to obtain actual values. Its nilpotent and
projective-incidence method builds on Tang's Appendix~A, while the
uniform seed, additive-core transfer and exact degree-eight and
degree-ten cases supply further constructions. The general even-degree
step uses the generic-nullity theorem of Brassil--Reichstein
\cite[Theorem~1.2]{BrassilReichstein2022} and Chen Lin, Chenhao Tang
and Enhan Zhao \cite[Theorem~1.2]{LinTangZhao2026}. The final
appendices give the exact finite checks supporting specified steps.

\part*{I. Odd standard polynomials}
\section{Linear slices and balanced endpoint transfer}\label{sec:slices}
Fix $A_1,\ldots,A_{m-1}\in\Mat_n(K)$ and set
\begin{equation}\label{eq:slice}
 L_A(X)=f(A_1,\ldots,A_{m-1},X).
\end{equation}
For multilinear $f$, the map $L_A$ is linear. Thus any linear
combination of outputs obtained by varying $X$ \emph{with the same fixed
tuple $A$} is again one value of $f$. No additive closure of the
unrestricted image is assumed. Simultaneous conjugation of all inputs
conjugates the output of $s_m$.

\begin{lemma}[Jordan-path criterion]\label{lem:jordanpath}
Suppose a slice image contains all diagonal matrices and the matrix
units
\[E_{i_1i_2},E_{i_2i_3},\ldots,E_{i_{n-1}i_n}\]
for some permutation $(i_1,\ldots,i_n)$ of the indices. Then
$\im f(\Mat_n(K))=\Mat_n(K)$, provided $f$ is equivariant under
simultaneous conjugation.
\end{lemma}
\begin{proof}
Within this one slice we can prescribe each diagonal entry and the
superdiagonal entries along the indicated path independently. A
Jordan matrix, including a direct sum of several blocks, has precisely
this form after permuting the coordinates, with its link entries
chosen to be $1$ or $0$. Every matrix over $K$ is similar to such a
matrix. Conjugating the entire input tuple concludes the proof.
\end{proof}

For later use let
\begin{equation}\label{eq:V2}
 V_h=\Span\{E_{i,i+h\bmod n}:0\le i<n\}.
\end{equation}
If $n$ is odd, the step-two edges in $V_2$ contain a Hamiltonian
cycle. If $n$ is even, the two parity cycles are linked by $E_{03}$;
an explicit Hamiltonian path is
\begin{equation}\label{eq:evenpath}
 2\longrightarrow4\longrightarrow\cdots\longrightarrow n-2
 \longrightarrow0\longrightarrow3\longrightarrow5
 \longrightarrow\cdots\longrightarrow n-1\longrightarrow1.
\end{equation}

Let $f$ be homogeneous multilinear of degree $m$ and let its final
variable be free. In the signed word expansion, denote by $P(A)$
the total fixed-input product coefficient when that variable is at
the \emph{last} position; let $Q(A)$ be its counterpart in the
\emph{first} position. Both depend polynomially on $A$.

\begin{theorem}[Balanced endpoint transfer]\label{thm:transfer}
Let $n=p+q$ with $|p-q|\le1$. Suppose there are fixed
$(m-1)$-tuples $A\in\Mat_p(K)^{m-1}$ and
$B\in\Mat_q(K)^{m-1}$ such that both slice images contain every
diagonal matrix in their respective blocks and $P(A)$ and $Q(A)$
are invertible. Then $\im f(\Mat_n(K))=\Mat_n(K)$, assuming
simultaneous-conjugacy equivariance.
\end{theorem}
\begin{proof}
Take the block-diagonal fixed tuple
\[
 C_j(t)=\begin{pmatrix}A_j&0\\0&tB_j\end{pmatrix}
 \qquad(1\le j<m).
\]
Writing the free input in four rectangular blocks yields four
independent block-preserving linear maps. On the two diagonal blocks
they are $L_A$ and $t^{m-1}L_B$. On a top-right input $X$,
setting $t=0$ kills terms in which a $B$ factor occurs on the
right of $X$, giving $X\mapsto P(A)X$. On a bottom-left input, the
limit is $X\mapsto XQ(A)$. These last maps are invertible for $t=0$.
The corresponding determinants are nonzero polynomials in $t$.
Choose a single $t\in K^\times$ avoiding their finitely many zero
sets. This choice is independent of the eventual target matrix.

At this fixed nonzero $t$, the slice covers all four mixed-block
entries and all diagonal entries in either square block. Since
$|p-q|\le1$, the bipartite graph of entries between the two blocks
has a Hamiltonian path alternating between the blocks. Therefore
Lemma~\ref{lem:jordanpath} applies.
\end{proof}

For $f=s_{2k+1}$ the two endpoint coefficients are
\begin{equation}\label{eq:endpoints}
 P(A)=Q(A)=s_{2k}(A_1,\ldots,A_{2k}).
\end{equation}
Indeed moving the free letter across the other $2k$ letters changes
no sign. The balanced hypothesis is essential: no Hamiltonian path
through a complete bipartite graph exists if its two parts differ in
size by more than one.

\section{Signed trails on a bidirected path}\label{sec:path}
Fix $k\ge1$, write $r=k+1$, and consider the ordered list
\begin{equation}\label{eq:pathlist}
 \mathcal A_k=(E_{01},E_{12},\ldots,E_{k-1,k},
 E_{k,k-1},E_{k-1,k-2},\ldots,E_{10}).
\end{equation}
Set
\[
\begin{aligned}
 c&=(1,\underbrace{2,\ldots,2}_{r-2},1)^t,\\
 H&=\diag(c)+cc^t,\\
 \mathcal E_r&=\diag(1,-1,1,\ldots,(-1)^{r-1}).
\end{aligned}
\]
For $r=2$ the middle list is empty. Every nonzero product of
matrix units is an oriented Euler trail through their directed
edges. This simple observation gives the exact slice coefficients.

\begin{lemma}[Endpoint and loops]\label{lem:pathloop}
For the tuple \eqref{eq:pathlist},
\begin{align}
 s_{2k}(\mathcal A_k)&=\diag((-1)^ic_i)_{0\le i<r},\label{eq:pathendpoint}\\
 s_{2k+1}(\mathcal A_k,\diag(x))&=\diag(H\mathcal E_r x).
 \label{eq:pathdiag}
\end{align}
Furthermore,
\begin{equation}\label{eq:pathdet}
 \det(H\mathcal E_r)=(-1)^{r(r-1)/2}2^{r-2}(2r-1)\ne0.
\end{equation}
\end{lemma}
\begin{proof}
The forward and reverse directed edges form a bidirected path. Every
closed Euler ordering is a cyclic rotation of the tour traversing
that path forward and returning. Vertex $i$ occurs as a starting
vertex $c_i$ times, and the resulting permutations have sign
$(-1)^i$. This proves \eqref{eq:pathendpoint}.

Insert a diagonal loop at vertex $j$. For a tour beginning at $i$,
there are $c_j+\delta_{ij}$ loop slots. Incorporating the loop into
the fixed ordering gives sign $(-1)^j$ for the added-loop
coefficient, so its contribution to output coordinate $i$ is
$c_i(c_j+\delta_{ij})(-1)^j$. This is exactly
$H\mathcal E_r$. Finally the determinant lemma yields
\[
 \det H=\left(\prod_{i=0}^{r-1}c_i\right)
 \left(1+c^t\diag(c)^{-1}c\right)
 =2^{r-2}(2r-1).
\]
Multiply by $\det\mathcal E_r=(-1)^{r(r-1)/2}$.
\end{proof}

\Needspace{18\baselineskip}
\begin{lemma}[Free chord on a bidirected path]\label{lem:pathchord}
Embed \eqref{eq:pathlist} on an interval of indices
$a,a+1,\ldots,a+k$, with $c_j$ as above. For $p\ne q$, the
image of the free matrix unit $E_{pq}$ is a multiple of $E_{pq}$.
The coefficient, before any external window weight, is
\begin{equation}\label{eq:pathchordtable}
\begin{array}{@{}ll@{}}
\text{position of $p,q$}&\text{coefficient}\\\midrule
p,q\text{ inside, same parity}&(-1)^{p-a}c_{p-a}c_{q-a}\\
p,q\text{ inside, opposite parity}&0\\
p\text{ inside},\ q\text{ outside}&(-1)^{p-a}c_{p-a}\\
q\text{ inside},\ p\text{ outside}&(-1)^{q-a}c_{q-a}\\
p,q\text{ outside}&0.
\end{array}
\end{equation}
\end{lemma}
\begin{proof}
The directed degree imbalance of the fixed bidirected path is zero.
Adding $p\to q$ forces every connected Euler trail to start at $p$
and end at $q$, so the output unit is $E_{pq}$. If the free edge
joins a path vertex to an external vertex, it must be first or
last; \eqref{eq:pathendpoint} gives the two one-sided coefficients.
If both endpoints lie outside, the selected edges are disconnected.

For two inside vertices, assume first $p<q$ and put $d=q-p$.
The tails outside $[p,q]$ must be traversed outward and back. Each
nonempty tail admits the independent attachment multiplicity
$c_{p-a}$ or $c_{q-a}$. After stripping these tails, an Euler trail
is specified by a turning vertex $j\in[p,q]$: go from $p$ to $j$
and back, take the chord $p\to q$, then go from $q$ to $j$ and
back. Moving $j$ one edge changes the permutation sign by $-1$.
The initial sign, by inversion counting in the displayed order of
\eqref{eq:pathlist}, is $(-1)^{q-a}$. Summing gives
\[
 (-1)^{q-a}c_{p-a}c_{q-a}
 \sum_{h=0}^{d}(-1)^h.
\]
The sum vanishes for odd $d$ and equals $1$ for even $d$, in
which case $p$ and $q$ have the same parity. For $p>q$, reversing
both the directed path and the fixed ordering contributes the
same sign $(-1)^k$ twice, giving the identical formula.
\end{proof}

\section{Positive path cores and large matrix sizes}\label{sec:large}
We next build one slice with invertible diagonal restriction and
invertible even-standard endpoint on every $\Mat_q(K)$ with $q\ge r$.
For $k=r-1$, define
\begin{equation}\label{eq:Vandermonde}
\begin{split}
 F_j&=\sum_{e=0}^{q-2}(e+1)^{j-1}E_{e,e+1},\\
 R_j&=\sum_{e=0}^{q-2}z_e(e+1)^{j-1}E_{e+1,e},
 \qquad z_e=(-1)^{\lfloor e/k\rfloor},
\end{split}
\end{equation}
for $1\le j\le k$, and order the fixed inputs as
$F_1,\ldots,F_k,R_k,\ldots,R_1$. Put
\[
 v_k=\prod_{j=1}^{k-1}j!,\qquad
 \iota_a:K^r\longrightarrow K^q,
 \quad(x_0,\ldots,x_{r-1})\longmapsto
 \sum_{i=0}^{r-1}x_i e_{a+i}.
\]

\begin{proposition}\label{prop:positive}
For every $q\ge r$, the diagonal restriction of the final-variable
slice associated with \eqref{eq:Vandermonde} is
\begin{equation}\label{eq:Gram}
 L_{\rm diag}=G\mathcal E_q,
 \qquad G=v_k^2\sum_{a=0}^{q-r}\iota_aH\iota_a^t.
\end{equation}
Its fixed-input even standard polynomial is
\begin{equation}\label{eq:largeendpoint}
 s_{2k}(F_1,\ldots,F_k,R_k,\ldots,R_1)
 =v_k^2\diag\left((-1)^i
 \sum_{\substack{0\le a\le q-r\\a\le i\le a+k}}
 c_{i-a}\right)_{0\le i<q}.
\end{equation}
Both matrices are invertible in characteristic zero.
\end{proposition}
\begin{proof}
Expand the fixed inputs into path units. Alternation among the
$F_j$ requires $k$ distinct forward edges, and alternation among
the $R_j$ requires $k$ distinct backward edges. A diagonal output
has zero net flow through every edge cut, so the forward and reverse
edge sets coincide. With only one extra diagonal loop, a disconnected
set of edges cannot support one Euler trail. Consequently the $k$
selected edges form a consecutive interval
$a,a+1,\ldots,a+k-1$.

On a fixed window, the coefficients from the forward and reverse
choices are two Vandermonde determinants, both of absolute algebraic
value $v_k$, before reverse signs. Their product is $v_k^2$.
Furthermore
\[
 \prod_{e=a}^{a+k-1}z_e=(-1)^a:
\]
the product is $+1$ at $a=0$, and shifting the window changes its
sign because $z_{a+k}/z_a=-1$. Multiplying the local formula
\eqref{eq:pathdiag} by this sign converts the local alternating
factor into the restriction of $\mathcal E_q$. Thus the windows
sum to \eqref{eq:Gram}; \eqref{eq:pathendpoint} similarly yields
\eqref{eq:largeendpoint}.

There is no assumption that $K$ is ordered. Instead regard these
integer matrices first over $\mathbb R$. The matrix
$H=\diag(c)+cc^t$ is positive definite, as is each of its
nonzero embedded quadratic forms. Every index $i$ occurs in at least
one interval $[a,a+k]$, so for $x\ne0$,
\[
 x^tGx=v_k^2\sum_{a=0}^{q-r}
 (\iota_a^tx)^tH(\iota_a^tx)>0.
\]
Therefore $\det G$ is a nonzero integer, and remains nonzero over
any field of characteristic zero. All diagonal entries in
\eqref{eq:largeendpoint} are nonzero signed positive integers.
\end{proof}

\begin{corollary}\label{cor:large}
For $r\ge2$ and $n\ge2r$, the image of $s_{2r-1}$ on
$\Mat_n(K)$ is all of $\Mat_n(K)$.
\end{corollary}
\begin{proof}
Let $p=\lceil n/2\rceil$, $q=\lfloor n/2\rfloor$.
Both are at least $r$; apply Proposition~\ref{prop:positive}
to each block and Theorem~\ref{thm:transfer}, using
\eqref{eq:endpoints}.
\end{proof}

\section{Minimal size: the cycle and its maximal-minor bivector}
\label{sec:minimal}
We treat the boundary
\begin{equation}\label{eq:minimalclaim}
 \im s_{2n-1}(\Mat_n(K))=\Mat_n(K),\qquad n\ge3.
\end{equation}
At this size, a single fixed slice need not serve all targets. The
construction instead adjusts the diagonal fixed inputs to each
prescribed Jordan matrix.

Fix the $n$ directed cycle edges
\[
 C_i=E_{i,i+1\bmod n},\qquad 0\le i<n,
\]
followed by $n-2$ diagonal matrices
$D_1,\ldots,D_{n-2}$ with diagonal row matrix
$D\in\Mat_{n-2,n}(K)$. Leave the final variable free and denote
the resulting slice by $L_D$. Let
\begin{equation}\label{eq:An}
 A_n=2\one\one^t-I_n.
\end{equation}
For a row-deleted coordinate matrix of rank $n-2$, the signed
maximal minors of $D$ form a skew matrix $B(D)$, described
explicitly in the following lemma. Its scale is fixed by the
order of the cycle inputs followed by the diagonal inputs.

\begin{lemma}[Cycle--bivector formulas]\label{lem:cycle}
For $0\le i<j<n$, let $D_{\widehat{i}\widehat{j}}$ be the
square submatrix obtained by removing columns $i,j$, and put
\begin{equation}\label{eq:gamma}
\begin{aligned}
 \gamma_{ij}&=(-1)^{n(n+1)/2+j-i-2},\\
 B_{ji}&=\gamma_{ij}\det D_{\widehat{i}\widehat{j}},
 \quad B_{ij}=-B_{ji},\quad B_{ii}=0.
\end{aligned}
\end{equation}
Then
\begin{equation}\label{eq:cycleDiag}
 L_D(\diag(x))=\diag(A_nB x).
\end{equation}
If $D$ has rank $n-2$, then $B$ has rank two and is decomposable.
Conversely, for independent $v,w\in K^n$, choosing the rows of
$D$ to span $\{x:v^tx=w^tx=0\}$ makes
$B$ a nonzero scalar multiple of $vw^t-wv^t$.

For $p\ne q$, let $I(p,q)$ be the inclusive forward cyclic
interval $q,q+1,\ldots,p$ and define
\begin{equation}\label{eq:ab}
 a_i=\begin{cases}
 1,&i=p\text{ or }i=q,\\2,&i\in I(p,q)\setminus\{p,q\},\\0,&i\notin I(p,q),
 \end{cases}\qquad
 b_i=\mathbf1_{i\notin I(p,q)}.
\end{equation}
The free chord $E_{pq}$ obeys the exact formula
\begin{equation}\label{eq:cycleChord}
 L_D(E_{pq})=-2(a^tBb)E_{pq}.
\end{equation}
In particular the original forward cycle edges have zero
coefficient.
\end{lemma}
\begin{proof}
All coefficients are alternating multilinear expressions in the
$n-2$ diagonal rows of $D$. It is therefore enough to calculate them
on coordinate diagonal inputs omitting a pair of vertices $i<j$;
expansion in minors then gives the stated general formulas.

For a free loop at $i$, the cycle and the $n-2$ fixed loops admit
closed Euler orderings obtained by rotating a cyclic ordering.
The vertex $j$ without a loop occurs once as a starting point and
all other vertices twice. Since the total word length is $2n-1$,
these rotations have the same sign. Place the diagonal loops in
increasing coordinate order after the cycle edges. Inversion
counting of the rotation and its loop insertions gives the common
factor
\[
 \gamma_{ij}=(-1)^{n(n+1)/2+j-i-2}.
\]
Thus the free loop at $i$ produces
$\gamma_{ij}\diag(2\one-e_j)$; the loop at $j$ gives the
negative corresponding expression with $i$ omitted. The two
columns are exactly $A_nB_{\cdot i}$ and $A_nB_{\cdot j}$.
This establishes \eqref{eq:cycleDiag} in the coordinate case.

For completeness we record the chord sign calculation in the same
minor normalization. An Euler trail containing the free chord
$p\to q$ has two possible edge skeletons before the fixed loops
are inserted: chord first followed by the forward cycle starting
at $q$, or the cycle from $p$ followed by the chord. A diagonal
loop strictly in the inclusive interval $I(p,q)$ has a forced
placement, while at either chord endpoint there are two placements.
Pairing the skeletons for a pair of omitted loops on the same side
of the interval cancels their signs. Moving an omitted loop
across an endpoint changes the inversion parity; the resulting
complete coefficient table, for the elementary bivector
$B_{ij}=-B_{ji}$, is
\begin{equation}\label{eq:chordsigns}
\begin{array}{@{}lc@{}}
\text{positions of omitted }i,j\ (i<j)&\text{coefficient of }B_{ij}\\\hline
 i,j\in I\text{ or }i,j\notin I&0\\
 i\text{ endpoint of }I,\ j\notin I&-2\\
 i\text{ strictly inside }I,\ j\notin I&-4\\
 i\notin I,\ j\text{ endpoint of }I&2\\
 i\notin I,\ j\text{ strictly inside }I&4.
\end{array}
\end{equation}
For example, take $n=4$, $p=2$, $q=0$ and omit the loop vertices
$i=0,j=3$, leaving $D_1=E_{11}$ and $D_2=E_{22}$.
Writing $X=E_{20}$ and keeping the input order
$C_0,C_1,C_2,C_3,D_1,D_2,X$, the four admissible Euler words and
permutation signs are
\[
\begin{array}{@{}lr@{}}
 C_2C_3C_0D_1C_1D_2X & -1\\
 XC_0D_1C_1D_2C_2C_3 & -1\\
 D_2C_2C_3C_0D_1C_1X & +1\\
 D_2XC_0D_1C_1C_2C_3 & -1.
\end{array}
\]
Their sum is $-2$, as in the endpoint/outside row of the table:
\eqref{eq:gamma} gives $\gamma_{03}=-1$, hence $B_{03}=1$.
Here is a parameter-uniform inversion check of the table, including
cyclic wraps. Write $S=\{0,\ldots,n-1\}\setminus\{i,j\}$, and
let $\ell=(p-q)\bmod n\in\{1,\ldots,n-1\}$. For the first
skeleton, with any endpoint loop placed at its earliest opportunity,
list the vertices in the order
\[
 \pi_{p,q}=(p,q,q+1,\ldots,p-1,p+1,\ldots,q-1),
\]
where each interval runs in the forward cyclic direction. Deleting
$i,j$ gives the loop ordering $\pi_{p,q}|_S$. Define
$\eta(p)=\eta(q)=n$ and
$\eta(v)=n-((v-q)\bmod n)$ otherwise, and write
$\operatorname{inv}$ for the number of numerical inversions.
The sign of this baseline Euler word is
\begin{equation}\label{eq:cycle-uniform-sign}
 \varepsilon=(-1)^{q(n-q)+\mathbf1_{p\in S}
       +\operatorname{inv}(\pi_{p,q}|_S)
       +\sum_{v\in S}\eta(v)}.
\end{equation}
The terms count, respectively, the rotation of the labelled cycle,
the last-labelled chord crossing the optional early $p$-loop,
inversions among the fixed-loop labels, and the cycle letters
following their respective loops. In particular, the endpoint
values of $\eta$ cannot be replaced by its interior formula.

To see that this absolute sign matches the cofactor convention, let
$z_v$ be the zero-based position of $v$ in $\pi_{p,q}$.
Moving $p$ from position $\ell$ to the front of the forward
rotation starting at $q$ yields
\[
 \operatorname{inv}(\pi_{p,q})\equiv q(n-q)+\ell,
 \qquad \sum_{v=0}^{n-1}\eta(v)=\frac{n(n+1)}2+\ell
 \pmod2.
\]
Deleting $v$ changes inversion parity by $z_v+v$; deleting both
$i<j$ requires the additional parity correction
$\mathbf1_{z_j<z_i}$. Substitution in
\eqref{eq:cycle-uniform-sign} gives
\[
 \frac{\varepsilon}{\gamma_{ij}}=
 (-1)^{1+h(i)+h(j)+\mathbf1_{z_j<z_i}},
 \quad h(v)=z_v+\eta(v)+\mathbf1_{v=p}\pmod2.
\]
The vertex order gives $h(v)\equiv n+1$ for $v\in I(p,q)$ and
$h(v)\equiv n$ outside. When $i,j$ lie on opposite sides,
$h(i)+h(j)\equiv1$, and all interval vertices precede the
outside vertices in $\pi_{p,q}$. Hence
\begin{equation}\label{eq:cycle-uniform-normalization}
 \varepsilon=\gamma_{ij}\operatorname{sgn}(a_i b_j-a_j b_i)
 \quad\text{whenever }a_i b_j-a_j b_i\ne0.
\end{equation}

For clarity, the remaining sign summation is finite independently
of $n$. If $A$ is the forward arc from $q$ to $p$ and $B$ its
complement, their edge lengths, including forced interior loops,
are $\alpha=2\ell-1-u$ and $\beta=2(n-\ell)-1-v$,
where $u,v$ count the omitted strict-interior loops in the
respective arcs. Let $P,Q\in\{0,1\}$ indicate the presence of
endpoint loops at $p,q$. In the first skeleton the parity changes
on moving a present $p$-loop or $q$-loop to its late placement
are $1+Q+\alpha$ and $\alpha+\beta$; in the second they are
$\alpha+\beta+Q$ and $\alpha+1$. The baseline skeletons differ
by $\beta(1+Q+\alpha)+Q+\alpha$ inversions.
Summing the at most $2^{P+Q+1}\le8$ possible words therefore
makes both same-side configurations cancel. In an opposite-side
configuration the signed sum divided by $\varepsilon$ is $2$
when the omitted interval vertex is an endpoint, and $4$
when it is a strict interior vertex. Together with
\eqref{eq:cycle-uniform-normalization}, the complete sum is
$2\gamma_{ij}(a_i b_j-a_j b_i)$, namely the table
\eqref{eq:chordsigns} with $B_{ij}=-\gamma_{ij}$.
Expansion in the maximal minors proves \eqref{eq:cycleChord}.
The degree imbalance of the chord forces output $E_{pq}$,
so no other matrix units occur.

The maximal-minor cofactor construction is the Pluecker bivector
of the two-dimensional annihilator of the row space of $D$.
For rank $n-2$, it is nonzero and decomposable, and its image is
that annihilator. In particular taking the row space to be
$\ker v^t\cap\ker w^t$ gives the claimed multiple of
$vw^t-wv^t$. No positivity or nonisotropy assumption on $v,w$
is involved.
\end{proof}

\begin{proposition}[Every Jordan target at the minimal size]
\label{prop:minJordan}
For each $n\ge3$ and every lower Jordan matrix $J\in\Mat_n(K)$,
there is a rank-$n-2$ choice of $D$ such that $J$ is contained in
the image of the single slice $L_D$, after a possible simultaneous
permutation of coordinates.
\end{proposition}
\begin{proof}
Write $d\in K^n$ for the diagonal of $J$ and $\tau=\sum_i d_i$.
The matrix $A_n$ is invertible: it acts by $2n-1$ on the line
$K\one$ and by $-1$ on its sum-zero complement. If $d\ne0$, put
\begin{equation}\label{eq:vdiag}
 v=A_n^{-1}d=-d+\frac{2\tau}{2n-1}\one.
\end{equation}
Select $w\notin Kv$ and take $B=vw^t-wv^t$, up to the nonzero
cofactor scalar. The range of $A_nB$ contains $d$. For a Jordan
link $E_{pq}$, the coefficient in \eqref{eq:cycleChord} is
proportional to
\begin{equation}\label{eq:linw}
 (a^tv)(w^tb)-(a^tw)(v^tb).
\end{equation}
If $a,b$ are independent, this is a nonzero linear functional of
$w$ precisely when $(a^tv,b^tv)\ne(0,0)$. The independence holds
for the nonforward edges used below: the interval is proper and
contains its two endpoints.

Here and below every Jordan block is represented with lower
subdiagonal links, and the coordinates may be permuted. For an
edge joining equal diagonal entries $\lambda$, formula
\eqref{eq:vdiag} yields
\begin{equation}\label{eq:intervalrelation}
 \frac12a^tv+b^tv=\lambda-\frac{\tau}{2n-1}.
\end{equation}
Indeed the left side assigns weight one to the two endpoints,
and weight one to every other coordinate; the sum is
$\sum_i v_i-(v_p+v_q)/2$. Substituting
$\sum_i v_i=\tau/(2n-1)$ and the common endpoint entries
$v_p=v_q=-\lambda+2\tau/(2n-1)$ gives the identity.

We consider all possible Jordan diagonals.
\smallskip

\noindent\emph{(i) $\tau\ne0$.}
List the Jordan blocks consecutively and use their ordinary lower
links $E_{i+1,i}$. Each relevant inclusive interval has just two
vertices. If $a^tv=b^tv=0$, the two equal endpoint entries of $v$
would sum to zero and so both vanish. The sum of the other entries
would vanish too, contradicting
$\sum_i v_i=\tau/(2n-1)\ne0$. For $n\ge3$ none of these lower
links is a forward cycle edge.

\smallskip
\noindent\emph{(ii) $\tau=0$ and $d\ne0$, with no nontrivial
zero-eigenvalue Jordan chain.}
Again use the consecutive Jordan ordering. Every required link
has nonzero eigenvalue $\lambda$; \eqref{eq:intervalrelation}
shows $a^tv$ and $b^tv$ cannot both vanish.

\smallskip
\noindent\emph{(iii) $\tau=0$, $d\ne0$, with at least one
nontrivial zero-eigenvalue Jordan chain.}
Two-color the vertices in every zero-eigenvalue Jordan chain as a
path, black and white, so each required zero link joins different
colors; assign any zero-eigenvalue singleton to either group.
Both colors occur. Select the leading endpoint of one nonzero-eigenvalue Jordan
chain as the single vertex in group $A$, with eigenvalue
$\mu\ne0$, and put the remaining nonzero-eigenvalue vertices
in group $B$, retaining their chain ordering. Since the total trace is zero, the sum of the
eigenvalues in $B$ is $-\mu$, so this group is nonempty. Arrange
the cyclic coordinates in the order
\begin{equation}\label{eq:blackwhite}
 \text{black zero vertices}\;|\; A\;|\;
 \text{white zero vertices}\;|\; B.
\end{equation}
A link with nonzero eigenvalue has nonzero functional by
\eqref{eq:intervalrelation}; its endpoints are not a forbidden
forward cycle edge. A zero Jordan link joins black to white,
and the inclusive cyclic interval from its source to its target
contains precisely one of $A,B$ among nonzero vertices. Thus
$\sum_{i\in I}d_i=\mu$ or $-\mu$. Now $v=-d$ and the zero
endpoints contribute nothing, so
$ a^tv=-2\sum_{i\in I}d_i=\mp2\mu\ne0$.
This deals simultaneously with several zero blocks and the
possible singleton vertices.

\smallskip
\noindent\emph{(iv) $d=0$.}
The target is nilpotent, possibly with several Jordan blocks.
Use the consecutive lower-Jordan ordering but choose $v=\one$
instead of \eqref{eq:vdiag}. The zero diagonal belongs to the
range of $A_nB$ automatically. For each required link the
inclusive interval has two vertices, so
$ a^tv=2$ and $b^tv=n-2$, and \eqref{eq:linw} is a
nonzero functional of $w$.

\smallskip
In each case we have finitely many nonzero linear functionals of
$w$ whose vanishing must be avoided, together with the proper
subspace $Kv$. Over the infinite field $K$, their union does not
exhaust $K^n$. Choose $w$ outside it, and select diagonal fixed
inputs realizing the decomposable bivector, as in
Lemma~\ref{lem:cycle}. Their coefficients on all required links
are nonzero. Solve \eqref{eq:cycleDiag} for an input $\diag(x)$
giving $d$ (with $x=0$ sufficient in (iv)), then add to \emph{that
same free input} the required $E_{pq}$ divided by their
coefficients in \eqref{eq:cycleChord}. The result is precisely the
permuted Jordan matrix $J$; the slice introduces no unwanted
matrix entries. Finally, conjugation gives the original target.
\end{proof}

\begin{corollary}\label{cor:minimal}
For every $n\ge3$, \eqref{eq:minimalclaim} holds.
\end{corollary}

\section{Intermediate matrix sizes: periodic path colors}
\label{sec:middle}
We now consider $r<n<2r$. Let $k=r-1$ and write
\begin{equation}\label{eq:kt}
 n=k+t+1,\qquad 1\le t\le k.
\end{equation}
First assume $k\ge3$. On the path with $n$ vertices, give
edge $e\to e+1$ the color $e\bmod k$. Put
\begin{equation}\label{eq:periodic}
\begin{split}
 F_j&=\sum_{\substack{0\le e<n-1\\e\equiv j\ (\mathrm{mod}\ k)}}
 E_{e,e+1},\\
 R_j&=\sum_{\substack{0\le e<n-1\\e\equiv j\ (\mathrm{mod}\ k)}}
 (-1)^{\lfloor e/k\rfloor}E_{e+1,e},\qquad 0\le j<k.
\end{split}
\end{equation}
Fix the inputs in the order
\begin{equation}\label{eq:periodorder}
 F_0,\ldots,F_{k-1},R_{k-1},\ldots,R_0
\end{equation}
and denote the final-variable slice by $L$.
Each term in the expansion chooses one forward and one reverse edge
of each of the $k$ colors. The colors $0,\ldots,t-1$ occur twice,
and the others once.

\subsection{Diagonal restriction and the connected-window rule}
A connected selection with the same forward and reverse edges must
be a window of $k$ consecutive edges. Every such window contains
each color once. The signs of the permutations of forward and
reverse colors cancel; the reversed-edge scalar signs have product
$(-1)^a$ when the window starts at $a$. By
Lemma~\ref{lem:pathloop},
\begin{equation}\label{eq:perdiag}
 L_{\rm diag}=\left(\sum_{a=0}^{t}
 \iota_a H\iota_a^t\right)\mathcal E_n.
\end{equation}
The Gram argument in Proposition~\ref{prop:positive} shows this
restriction is invertible: $H$ is positive definite over the reals,
and every coordinate belongs to one of the $t+1$ windows. Its
determinant is a nonzero integer.

We will use ordinary offsets as well as cyclic ones. For
$|h|<n$, put
\[
 \mathcal U_h=\Span\{E_{i,i+h}:0\le i,i+h<n\}.
\]
Because \eqref{eq:periodorder} contributes $k$ forward and $k$
backward unit steps, inserting $E_{u,u+h}$ gives output in
$\mathcal U_h$. Thus $L$ preserves \emph{ordinary} offsets,
and consequently respects their summands inside a cyclic offset.

A flow comparison across each path cut shows that, if the
coefficient from $E_{u,u+h}$ to $E_{i,i+h}$ is nonzero, the
multisets of colors on the oriented intervals beginning at $u$
and at $i$ of length $|h|$ must coincide. Removing complete cycles
of $k$ colors, a remaining nonzero proper color interval has a
unique start residue. Hence
\begin{equation}\label{eq:colorresidue}
 h\not\equiv0\pmod k\quad\Longrightarrow\quad i\equiv u\pmod k.
\end{equation}
This restriction is particularly strong since $n\le2k+1$:
for offsets two and three there are at most two positions in a
residue class.

For a coefficient along the diagonal of an ordinary-offset
restriction, input and output have the same endpoints, hence
forward and reverse selected edge sets agree. Encode the choice
of edge from duplicated color $j<t$ by $\epsilon_j=0$ for its
first occurrence and $\epsilon_j=1$ for its second. The edge
indicator along the path is
\begin{equation}\label{eq:binary}
 (1-\epsilon_0,\ldots,1-\epsilon_{t-1},
 \underbrace{1,\ldots,1}_{k-t},
 \epsilon_0,\ldots,\epsilon_{t-1}).
\end{equation}
When $t<k$, this pattern has one plus twice the number of
$0\to1$ rises of $\epsilon$ as its number of connected
components; the only permissible connected choices are the
ordinary windows. A single free chord cannot make one Euler
trail out of three or more components. When $t=k$, there is in
addition a two-component selection
$\epsilon=0^a1^{k-a}$ for each $1\le a<k$:
it takes one edge interval at the far left and another at the far
right, separated by a gap of $k$ edges. Such a selection cannot
be joined by a chord of length two if $k\ge3$, nor by a chord
of length three if $k\ge4$. We retain it for the long wrapping
offsets below.

\subsection{The nonwrapping offset-two restriction}
Consider $L|_{\mathcal U_2}$. By \eqref{eq:colorresidue},
its matrix splits into singleton blocks and $2\times2$ blocks
on the positions $i,i+k$ with $0\le i\le t-2$.
For a diagonal coefficient, the connected-window rule and
Lemma~\ref{lem:pathchord} show that all contributions have sign
$(-1)^i$, and their unsigned contributions are positive integers.
Some window contains both endpoints of $E_{i,i+2}$. For $k\ge3$
the multiplicity product in that window is at least two, so
\begin{equation}\label{eq:diagbound}
 |(L|_{\mathcal U_2})_{ii}|\ge2.
\end{equation}
For a pair $i,i+k$, the off-diagonal coefficient has absolute
value one. For example, to carry the free input
$E_{i+k,i+k+2}$ into output $E_{i,i+2}$, the cut equations force
the unmatched forward edges at $i,i+1$ and reverse edges at
$i+k,i+k+1$. Connectivity forces every edge in between. There
is then one Euler trail: forward from $i$ to $i+k$, take the
free chord, and return backward to $i+2$. Reversing this
construction gives the reverse coefficient, again of magnitude
one. The two-component exceptional selections cannot intervene
because their gap exceeds two.

Thus a paired restriction block has diagonal product of absolute
value at least four and off-diagonal product of absolute value
one; its determinant is nonzero. Every singleton is nonzero by
\eqref{eq:diagbound}. Hence
\begin{equation}\label{eq:U2}
 L(\mathcal U_2)=\mathcal U_2.
\end{equation}

\subsection{The wrapping offset-two restriction}
The remaining two units in $V_2$ have ordinary offset $2-n$.
It is convenient to compute the opposite offset $n-2$ and use
an anti-involution. Choose a diagonal $Z$ with
\begin{equation}\label{eq:Z}
 \frac{Z_{e+1,e+1}}{Z_{ee}}=(-1)^{\lfloor e/k\rfloor}
 \qquad(0\le e<n-1),
\end{equation}
and define $\theta(X)=ZX^tZ^{-1}$.
The map $\theta$ reverses products and interchanges the forward
and reverse colored inputs. Reversing the word of length $2k+1$
contributes $(-1)^k$, and reversing the $2k$ fixed inputs
contributes the same sign. Consequently
\begin{equation}\label{eq:theta}
 L\theta=\theta L,
\end{equation}
so opposite ordinary-offset restrictions are conjugate and have
the same invertibility status.

There are two positions on $\mathcal U_{n-2}$. Its restriction
is given by the following complete table (the displayed ordered
pairs are diagonal entries):
\begin{equation}\label{eq:wrapping}
\begin{array}{@{}lll@{}}
 \text{case}&\text{diagonal entries}&\text{off-diagonal entries}\\\hline
 2\le t<k,\ n\text{ even}&(4,-4)&0\\
 2\le t<k,\ n\text{ odd}&(-2,-2)&0\\
 t=1,\ k\text{ even}&(3,-3)&\text{each of magnitude}\le1\\
 t=1,\ k\text{ odd}&(-2,-2)&\text{each of magnitude}\le1\\
 t=k&(-(2k-1),-(2k-1))&0.
\end{array}
\end{equation}
We give the signed counts because the last row does not follow
from the ordinary connected-window rule alone.

For $2\le t<k$, color conservation prevents mixing between the
two positions. Applying \eqref{eq:pathchordtable} to the
connected windows gives respectively
$1+3(-1)^{n-2}$ and $-3+(-1)^{n-1}$.
For $t=1$, the two window contributions give
$\mathbf1_{2\mid k}+2(-1)^k$ and
$-2-\mathbf1_{2\mid k}$. A possible off-diagonal trail has a
turning point at one of $k$ positions and the signs alternate,
so its sum is $0$ or $\pm1$. In either case the determinant is
nonzero.

If $t=k$, the connected windows contribute $-2$ to each
of the two diagonal coefficients. For the extra selections
$\epsilon=0^a1^{k-a}$, the left and right bidirected components
have lengths $a$ and $k-a$ and the reverse-color sign is
$(-1)^{k-a}$. For the long chord $E_{0,2k-1}$, the local
starting sign at the left endpoint is positive, and the sign of
the right component at its penultimate vertex is
$(-1)^{k-a-1}$. Thus every two-end contribution has negative
sign after the reverse-color sign is included. Its magnitude
is two except for $a=k-1$, where it is one. For the other long
chord $E_{1,2k}$, the left local sign is negative and the right
endpoint sign is $(-1)^{k-a}$; again all contributions are
negative, with the unique magnitude-one case $a=1$. More explicitly, for $b=k-a$ the contributions of one
selection to the respective diagonal positions are
\begin{equation}\label{eq:uniform-wrap-selections}
 -(2-\mathbf1_{a=k-1}),\qquad
 -(2-\mathbf1_{a=1}),\quad 1\le a\le k-1.
\end{equation}
The reverse-color sign $(-1)^b$ is included in both expressions.
Summing the $k-1$ choices gives $-(2k-3)$ in each position.
Including the connected-window contributions $-2$ gives
$-(2k-1)$; off-diagonal entries vanish by
\eqref{eq:colorresidue}, since the two long-chord positions
have different residues modulo $k\ge3$. This proves the
last row of \eqref{eq:wrapping} uniformly in $k$, including
its signs.

Every restriction in \eqref{eq:wrapping} is invertible. By
\eqref{eq:theta}, so is the restriction on $\mathcal U_{2-n}$.
Combining it with \eqref{eq:U2} yields
\begin{equation}\label{eq:V2full}
 L(V_2)=V_2.
\end{equation}
For odd $n$, \eqref{eq:perdiag} and \eqref{eq:V2full}
now imply surjectivity by Lemma~\ref{lem:jordanpath}.

\subsection{The additional edge for even sizes}
Assume $k\ge4$. For the ordinary-offset-three restriction,
\eqref{eq:colorresidue} isolates the residue-zero block. Its
index set is $\{0\}$ if $t\le2$, and $\{0,k\}$ if $t\ge3$.
The corresponding matrix is
\begin{equation}\label{eq:grade3}
 \begin{cases}
 [-2],&t=1,\\{}[-4],&t=2,\\
 \begin{pmatrix}-5&\varepsilon\\\eta&5(-1)^k\end{pmatrix},
  &3\le t\le k,
 \end{cases}
 \qquad |\varepsilon|=|\eta|=1.
\end{equation}
For the row at $0$, the windows starting at $1,2,3$ contribute
$-2,-2,-1$ when present; the window at zero contributes zero,
since its endpoints have opposite parity. For the row at $k$,
the first three windows contribute $1,2,2$, with common sign
$(-1)^k$. The paired off-diagonal coefficients are the
unique Euler-trail values $\pm1$, as for ordinary offset two;
the two-end selections have too long a gap. This proves
\eqref{eq:grade3}. Every block is invertible: in the last
case its diagonal product has absolute value $25$, while its
off-diagonal product has absolute value $1$. Hence the slice
contains $E_{03}$, by choosing the appropriate free input
inside the isolated residue-zero block. For even $n$, the
Hamiltonian path \eqref{eq:evenpath} completes the argument.

\begin{proposition}\label{prop:middle}
For $r\ge5$ and $r<n<2r$, the image of $s_{2r-1}$ is
$\Mat_n(K)$. The same argument gives this conclusion for
$(r,n)=(4,5)$ and $(4,7)$.
\end{proposition}
\begin{proof}
For $k\ge4$ the diagonal and cyclic offset-two subspaces are
in the same fixed slice by \eqref{eq:perdiag} and
\eqref{eq:V2full}; for even $n$ it also contains $E_{03}$
by \eqref{eq:grade3}. Apply Lemma~\ref{lem:jordanpath}.
If $k=3$, the relevant intermediate odd sizes are $n=5,7$;
step two is Hamiltonian and the same diagonal and offset-two
arguments apply. The even intermediate size $n=6$ is treated
separately below.
\end{proof}

\section{Exceptional small matrix orders}\label{sec:seams}
The preceding uniform proofs leave $s_5$ on $\Mat_4,\Mat_5$
and $s_7$ on $\Mat_6$. We now give concrete slices or a block
construction for each. In the two fixed-slice cases, the complete
restriction matrices below are exact integer certificates; the
recurrences for regenerating every entry are given in
Appendix~\ref{app:verification}.

\subsection{\texorpdfstring{$s_5$ on $\Mat_5$}{s5 on M5}}
Use \eqref{eq:periodic} with $k=t=2$ and the fixed-input order
\eqref{eq:periodorder}. This case is not covered by the previous
proof, which assumed $k\ge3$. The restrictions to the diagonal
and cyclic step-two spaces, in the natural unit bases, are
\begin{equation}\label{eq:s5mat}
L_0=\begin{pmatrix}
2&-2&1&0&0\\
2&-8&4&-1&0\\
1&-4&10&-4&1\\
0&-1&4&-8&2\\
0&0&1&-2&2
\end{pmatrix},\qquad
L_2=\begin{pmatrix}
4&-1&1&0&0\\
-1&-4&-1&0&0\\
1&-1&4&0&0\\
0&0&0&-3&0\\
0&0&0&0&-3
\end{pmatrix}.
\end{equation}
Direct exact evaluation gives
\begin{equation}\label{eq:s5dets}
 \det L_0=900,\qquad\det L_2=-594.
\end{equation}
Both restrictions are invertible. Since step two is Hamiltonian
on five vertices, Lemma~\ref{lem:jordanpath} proves
\begin{equation}\label{eq:s5m5}
 \im s_5(\Mat_5(K))=\Mat_5(K).
\end{equation}

\subsection{\texorpdfstring{$s_7$ on $\Mat_6$}{s7 on M6}}
Take the Vandermonde core \eqref{eq:Vandermonde} with
$r=4$ and $q=6$. Its diagonal restriction has determinant
\begin{equation}\label{eq:s7d0}
 \det L_0=-31\,911\,936.
\end{equation}
The cyclic offset-two restriction is
\begin{equation}\label{eq:s7mat2}
 L_2=\begin{pmatrix}
 -16&-24&24&-4&0&0\\
 24&-132&144&-24&0&0\\
 -24&144&60&24&0&0\\
 4&-24&24&16&0&0\\
 0&0&0&0&-416&-36\\
 0&0&0&0&-36&-376
 \end{pmatrix}.
\end{equation}
The exact determinant of this restriction is
\[\det L_2=1\,838\,804\,889\,600.\]
For the parity-linking step the cyclic offset-three restriction is
\begin{equation}\label{eq:s7mat3}
 L_3=\begin{pmatrix}
164&44&0&0&0&0\\
44&432&100&0&0&0\\
0&100&52&0&0&0\\
0&0&0&164&-44&0\\
0&0&0&-44&432&-100\\
0&0&0&0&-100&52
\end{pmatrix},
\quad\det L_3=3\,776\,896\,843\,776.
\end{equation}
All three restrictions belong to the \emph{same} final-variable
slice, and all are invertible. Thus it contains every diagonal,
all of $V_2$, and in particular $E_{03}\in V_3$. The even-size
Hamiltonian path \eqref{eq:evenpath} gives
\begin{equation}\label{eq:s7m6}
 \im s_7(\Mat_6(K))=\Mat_6(K).
\end{equation}

\subsection{\texorpdfstring{$s_5$ on $\Mat_4$}{s5 on M4}: an explicit three-dimensional core}
We use an elementary $3\times3$ construction. Put
\begin{equation}\label{eq:s5core}
 A=E_{01},\quad B=E_{12},\quad C=E_{20},\quad
 D=\diag(-1,0,1)
\end{equation}
inside $\Mat_3(K)$. Direct expansion of the alternator yields
\begin{equation}\label{eq:s4endpoint}
 s_4(A,B,C,D)=\diag(-1,2,-1),
\end{equation}
and the slice image is the five-dimensional subspace
\begin{equation}\label{eq:W}
 W=\left\{
 \begin{pmatrix}\alpha&0&z\\u&\beta&0\\0&v&\alpha\end{pmatrix}
 :\alpha,\beta,u,v,z\in K\right\}.
\end{equation}
Indeed the specified matrix is obtained from the single free input
\begin{equation}\label{eq:Wpre}
 X=\diag\left(\frac{3\beta-4\alpha}{5},
 \frac{2\beta-\alpha}{5},0\right)
 +\frac u2E_{10}+\frac v2E_{21}-\frac z4E_{02}.
\end{equation}
Equations \eqref{eq:s4endpoint} and \eqref{eq:Wpre} follow by
direct evaluation.

\begin{lemma}\label{lem:Wclasses}
Every similarity class in $\Mat_3(K)$ meets $W$.
\end{lemma}
\begin{proof}
For a cyclic target with characteristic polynomial
$x^3-tx^2+qx-\delta$, choose $\alpha\in K$ solving
\[
 2t\alpha-3\alpha^2=q,
 \qquad \beta=t-2\alpha,
 \qquad u=v=1,\qquad z=\delta-\alpha^2\beta.
\]
The matrix in \eqref{eq:W} has the required characteristic
polynomial. Moreover
$\det[e_0,Me_0,M^2e_0]=1$, so it is cyclic and the
characteristic polynomial identifies its cyclic similarity class.
For noncyclic nonscalar classes the elementary representatives
$\diag(c,d,c)$ and $cI+E_{10}$, as appropriate to the Jordan
partition, lie in $W$. Scalars plainly do also.
\end{proof}

Embed \eqref{eq:s5core} into the upper-left $3\times3$ block of
$\Mat_4(K)$, with zero lower-right singleton. Because the
even-standard endpoint \eqref{eq:s4endpoint} is invertible, the
two off-diagonal rectangular maps are arbitrary, and the
upper-left block ranges over $W$. The lower-right output is
always zero. Thus this one slice has image
\begin{equation}\label{eq:3plus1}
 \left\{\begin{pmatrix}Y&b\\c&0\end{pmatrix}:
 Y\in W,\ b\in K^{3\times1},\ c\in K^{1\times3}\right\}.
\end{equation}

Let $T\in\Mat_4(K)$ be nonscalar. There exists $v$ such that
$v,Tv$ are independent; otherwise every vector is an eigenvector
and $T$ would be scalar. Choose $\phi\in(K^4)^*$ with
$\phi(v)=1$, $\phi(Tv)=0$. In a basis adapted to
$K^4=\ker\phi\oplus Kv$, the bottom-right coefficient of $T$
is zero. By Lemma~\ref{lem:Wclasses}, a similarity change on
$\ker\phi$ makes its $3\times3$ compression a member of $W$.
The other two blocks are unrestricted in \eqref{eq:3plus1}.
Hence every nonscalar $T$ is a single $s_5$ value, up to
conjugation.

To obtain scalars we need a short general observation.
\begin{lemma}[Traceful multilinear values]\label{lem:scalar}
Let $f$ be multilinear on $\Mat_n(K)$, with $K$ algebraically
closed of characteristic zero. If $f$ has an evaluation with
nonzero trace, then some evaluation of $f$ is a nonzero scalar
matrix. Consequently every scalar matrix is in its actual image.
\end{lemma}
\begin{proof}
Fix a primitive $n$th root of unity and the usual clock-and-shift
matrices $U,V$ with $UV=\zeta VU$. The $n^2$ matrices
$U^aV^b$, $a,b\bmod n$, form a basis of $\Mat_n(K)$ and are
homogeneous for the $\mathbb Z_n^2$ grading. Every product of
a fixed tuple of homogeneous basis elements lies in the same
graded component. By multilinearity, expanding a traceful
evaluation in this basis gives a homogeneous basis tuple on which
$f$ has nonzero trace. Every nonzero graded component has trace
zero, whereas the degree-zero component is $KI$. Thus this
particular value is a nonzero scalar. Rescaling one input gives
any scalar.
\end{proof}

The core \eqref{eq:s5core} is traceful: for example,
\eqref{eq:W} contains $\diag(1,0,1)$, which is a single
$s_5$ value on $\Mat_3(K)$. Embedding the same evaluation in
$\Mat_4(K)$ preserves its nonzero trace. Lemma~\ref{lem:scalar}
therefore gives all scalar targets, completing
\begin{equation}\label{eq:s5m4}
 \im s_5(\Mat_4(K))=\Mat_4(K).
\end{equation}

\section{Proof of the odd-standard assertion}\label{sec:assembly}
\begin{proof}[Proof of the odd assertion of Theorem~\ref{thm:main}]
If $n<r$, then $2r-1\ge2n+1$. By Amitsur--Levitzki,
$s_{2n}$ is an identity on $\Mat_n(K)$; the standard recursion
expresses every higher alternator through lower ones, proving
$s_{2r-1}=0$.

For $r=1$, the claim is $s_1(X)=X$. For $r=2$ and $n=1$,
the alternating cubic vanishes. For $n\ge2$, the published
cubic image-linearity result \cite[Theorem~3.5]{Vitas2026}
implies that $\im s_3$ is a vector subspace. The substitution
$s_3(I,A,B)=[A,B]$ gives every traceless matrix by the classical
commutator theorem \cite{AlbertMuckenhoupt1957}, and a traceful
$2\times2$ evaluation,
embedded into $\Mat_n(K)$, gives a value outside that subspace.
For example, $s_3(E_{00},E_{01},E_{10})=\diag(2,1)$
on $\Mat_2(K)$. Since the traceless subspace has codimension one, cubic
image-linearity now gives the whole $\Mat_n(K)$.

Suppose $r\ge3$. The case $n=r$ follows from
Corollary~\ref{cor:minimal}. The range $n\ge2r$ follows from
Corollary~\ref{cor:large}. For $r<n<2r$, when $r\ge5$ apply
Proposition~\ref{prop:middle}; when $r=4$, it covers $n=5,7$
and \eqref{eq:s7m6} covers $n=6$. When $r=3$, the sizes
$n=4,5$ are \eqref{eq:s5m4} and \eqref{eq:s5m5}.
The cases are partitioned as follows:
\begin{center}
\small
\begin{tabular}{@{}ll@{}}
\toprule
Degree parameter & Surjective sizes and argument \\
\midrule
$r=1$ & all $n\ge1$; identity map\\
$r=2$ & all $n\ge2$; cubic theorem and direct values\\
$r=3$ & $n=3$ cycle, $n=4,5$ cases, $n\ge6$ transfer\\
$r=4$ & $n=4$ cycle; $n=5,7$ periodic; $n=6$ case; $n\ge8$ transfer\\
$r\ge5$ & $n=r$ cycle; $r<n<2r$ periodic; $n\ge2r$ transfer\\
\bottomrule
\end{tabular}
\end{center}
Together with the identity branch $n<r$, this exhausts every
pair of positive integers $(r,n)$.
\end{proof}

\begin{remark}
The large and periodic ranges use a fixed linear slice, whereas the
minimal-size construction chooses a slice for the prescribed Jordan target.
\end{remark}


\part*{II. Even standard polynomials}
\section{Actual images from an attained rank-one nilpotent}\label{sec:even-prelim}
The remainder of the article deals with $s_{2r}$. Throughout this part,
its trace-zero property follows directly from cyclic alternation:
for any word of length $2r$, moving the first letter to the last
preserves the trace but changes the permutation sign by
$(-1)^{2r-1}=-1$. Consequently
\begin{equation}\label{eq:even-trace}
 \tr s_{2r}(A_1,\ldots,A_{2r})=0,
 \qquad \im s_{2r}(\Mat_n)\subseteq\Sl_n.
\end{equation}
The identity $s_{2n}=0$ on $\Mat_n$ implies the entire zero branch
for $n\le r$. Moreover $s_{2r}(A_1,\ldots,A_{2r-1},I_n)=0$.
For nonzero images the following local-to-global mechanism is essential.
It is the smooth-map formulation of the nilpotent argument in
\cite[Lemma~2]{Tang2026FiveSix}; the elementary degeneration is
included to make the actual-value conclusion explicit.

\begin{lemma}[Attained-nilpotent actual-image criterion]
\label{lem:nilpotent-open}
Let $F:\mathbb A^m_K\to\Sl_n(K)$ be a polynomial map whose image
contains zero and is invariant under similarities and multiplication
by nonzero scalars. Suppose that for some $x$ the value $R=F(x)$
is nonzero, rank one and square zero, and
\(dF_x:K^m\to\Sl_n(K)\) is surjective. Then
$\im F=\Sl_n(K)$.
\end{lemma}
\begin{proof}
Both affine spaces are smooth. Surjectivity of the differential
makes $F$ smooth at $x$, and smoothness is open; the restriction to
a suitable neighborhood maps onto a Zariski-open subset containing
$R$. In particular, the image contains an open neighborhood $U$ of
$R$ in $\Sl_n$.

Let $T\in\Sl_n$ be nonzero. It is nonscalar in characteristic zero.
A nonscalar linear transformation has a vector $v$ for which
$v,Tv$ are linearly independent: otherwise every line is invariant,
forcing the transformation to be scalar. Thus after a similarity
$T_{01}\ne0$. For $t\ne0$, write
\begin{equation}\label{eq:nilpotent-degeneration}
 D(t)=\diag(t^{-1},t,1,\ldots,1),\qquad
 T(t)=t^2D(t)TD(t)^{-1}.
\end{equation}
All entries of $T(t)$ have nonnegative powers of $t$; its value at
$t=0$ is $T_{01}E_{01}$. Every nonzero rank-one square-zero matrix
is similar, up to a nonzero scalar, to $R$. By similarity and
scaling invariance there is therefore an open neighborhood of
$T_{01}E_{01}$ contained in the image. Its preimage under the
affine-line morphism $t\mapsto T(t)$ contains zero and hence some
nonzero $t$. Thus $T(t)$ is an actual value. Undoing the
similarity and scalar $t^2$ shows that $T$ is an actual value.
The zero target is already in the image.
\end{proof}

A standard polynomial is equivariant for simultaneous conjugation;
scaling one input scales its output, and setting an input to zero
gives zero. Thus every standard evaluation map satisfies the
invariance assumptions of Lemma~\ref{lem:nilpotent-open}.

\section{A uniform minimal-size seed}\label{sec:even-seed}
The first matrix order at which $s_{2r}$ need not vanish is
$n=r+1$. The following construction treats every such order with
one matrix-unit formula, including the full differential required
by Lemma~\ref{lem:nilpotent-open}.

\begin{theorem}[Uniform even-standard seed]\label{thm:even-seed}
For every $n\ge2$,
\begin{equation}\label{eq:even-seed}
 \im s_{2n-2}(\Mat_n(K))=\Sl_n(K).
\end{equation}
More precisely, the evaluation constructed below has a nonzero
rank-one square-zero output and surjective full differential.
\end{theorem}
\begin{proof}
Write $p=n-2$ and index matrix units by $0,\ldots,n-1$.
For $n\ge3$, use the ordered $2p+2$ inputs
\begin{equation}\label{eq:even-tuple}
\begin{aligned}
 \mathcal A=(&L=E_{00},\ X=E_{01},\\[-2pt]
 &S_2=E_{02},T_2=E_{20},\ldots,\\[-2pt]
 &S_{n-1}=E_{0,n-1},T_{n-1}=E_{n-1,0}).
\end{aligned}
\end{equation}
The signed matrix-unit trail calculation gives
\begin{equation}\label{eq:even-seed-evals}
\begin{aligned}
 s_{2n-2}(\mathcal A)&=(p+1)!E_{01},\\
 \partial_{X\gets E_{10}}s_{2n-2}(\mathcal A)&=-(p+1)!E_{10},\\
 \partial_{X\gets E_{12}}s_{2n-2}(\mathcal A)&=-p!E_{12},\\
 \partial_{S_2\gets E_{10}}s_{2n-2}(\mathcal A)&=-(p+1)!E_{20},\\
 \partial_{S_2\gets E_{12}}s_{2n-2}(\mathcal A)&=p!(E_{11}-E_{22}).
\end{aligned}
\end{equation}
The notation $\partial_{A_j\gets V}$ denotes the directional
derivative in the $j$th argument; by multilinearity it is exactly
the value with that argument replaced by $V$.

Here is a signed-trail derivation valid for all $n$. A nonzero
matrix-unit word is precisely an oriented Euler trail on the
labelled directed graph. In the unmodified tuple, each trail starts
at $0$, traverses the loop $L$ and the $p$ excursions
$S_jT_j$ in any order, and ends with $X$. The excursions have
even length and commute in permutation parity; moving $X$ to the
end crosses $2p$ labels. These $(p+1)!$ trails all have positive
sign. With $X=E_{10}$ the trail instead begins with $X$, making
its sign negative. With $X=E_{12}$ it starts with $X,T_2$,
traverses the remaining root blocks and ends with $S_2$;
there are $p!$ such trails and the word has odd inversion parity.
With $S_2=E_{10}$ the trail begins with $T_2$ and traverses the
root loop, the new excursion $XS_2$, and the remaining
$p-1$ excursions; the resulting $(p+1)!$ trails are negative.

Finally, for $S_2=E_{12}$, there is a directed triangle
$X,S_2,T_2$ together with a loop and the other two-edge
excursions. Trails starting at $0$ or another leaf cancel by
exchanging the odd-length loop and triangle. Starting at $1$
leaves $p!$ positively signed trails, whereas starting at $2$
leaves $p!$ negatively signed trails. This proves the last
formula in \eqref{eq:even-seed-evals}, including its diagonal
signs.

Let $\mathcal D$ be the image of the \emph{full} differential of
$s_{2n-2}$ at $\mathcal A$. It lies in $\Sl_n$ by
\eqref{eq:even-trace}. Equivariance supplies every commutator
$[Z,E_{01}]\in\mathcal D$ (the scalar $(p+1)!$ is nonzero).
Consider a change of basis $g=\diag(1,1,G)$,
$G\in\mathrm{GL}_{p}(K)$, acting on the leaf coordinates.
It fixes $L$ and $X$ and sends the $p$ outgoing and $p$ incoming
leaf arguments to mutually contragredient linear combinations
of their original arguments. The combined transformation on
these $2p$ frozen slots has determinant one. The alternating
multilinear map, and hence its differential image at the tuple,
is therefore invariant under conjugation by $g$.

The orbits of $E_{12}$ and $E_{20}$ in
\eqref{eq:even-seed-evals} supply all $E_{1j}$ and $E_{i0}$
for $i,j\ge2$. The orbit of $E_{11}-E_{22}$ spans
\begin{equation}\label{eq:even-seed-block}
 \bigl\{(\tr B)E_{11}-B:
 B\text{ is supported on the leaf block}\bigr\},
\end{equation}
because the conjugates of a rank-one diagonal idempotent span
its full matrix block: rank-one idempotents include $E_{ii}$
and $E_{ii}+E_{ji}$ for distinct leaf indices.
Together with $E_{10}$ and $[Z,E_{01}]$, these directions give
every off-diagonal matrix unit and every diagonal difference.
Thus $\mathcal D=\Sl_n$, as required.

When $n=2$, take $(E_{00},E_{01})$.
Its commutator value is $E_{01}$; replacing $E_{01}$ by $E_{10}$
gives $-E_{10}$ and the conjugation direction gives
$E_{00}-E_{11}$. The full differential is $\Sl_2$.
In each order the output is a nonzero rank-one square-zero
matrix, so Lemma~\ref{lem:nilpotent-open} proves
\eqref{eq:even-seed}.
\end{proof}

\begin{remark}[The size-four seed in Tang's construction]\label{rem:tang-seed}
At $n=4$, the tuple in \eqref{eq:even-tuple} is
$(E_{00},E_{01},E_{02},E_{20},E_{03},E_{30})$.
Qiyue Tang's Appendix~A uses
$(4E_{11},E_{12},E_{13},E_{14},E_{31},E_{41})$ in one-based
coordinates \cite[Proposition~11]{Tang2026FiveSix}.
The two coincide after relabelling the indices, interchanging
the fourth and fifth slots, and multiplying the first slot by
$4$; alternating multilinearity accounts for the sign and
scale. Thus the size-four instance agrees, up to these changes,
with Tang's matrix-unit construction.
\end{remark}

\section{Admissible cores and an additive transfer}\label{sec:even-cores}
The seed can be transferred to larger orders by a normal-map
criterion that keeps track of the \emph{full differential}, rather
than merely of a frozen generic kernel. Fix $r\ge1$ and write
\begin{equation}\label{eq:even-rk}
 k=2r-1.
\end{equation}
For $B=(B_1,\ldots,B_k)\in\Mat_N(K)^k$ and
$u\in K^N$, define
\begin{equation}\label{eq:core-definitions}
\begin{aligned}
 S&=s_k(B_1,\ldots,B_k),\\
 L(P)&=s_{k+1}(B_1,\ldots,B_k,P),\\
 C_i(P)&=s_k(B_1,\ldots,B_{i-1},P,B_{i+1},\ldots,B_k),\\
 \Gam_{B,u}(P)&=\bigl(L(P),C_1(P)u,\ldots,C_k(P)u\bigr).
\end{aligned}
\end{equation}
Thus $\Gam_{B,u}:\Mat_N(K)\to\Mat_N(K)\oplus(K^N)^k$
is a linear map in $P$; the frozen data $(B,u)$ vary
polynomially. Call the size $N$ \emph{admissible for $r$} if
there exists one such tuple with
\begin{equation}\label{eq:core-admissible}
 \det S\ne0,\qquad \ker\Gam_{B,u}=\{0\}.
\end{equation}
Both conditions are polynomial nonvanishing conditions in
unrestricted parameters. A witness therefore defines a nonempty
Zariski-open subset of the core parameter space.

\begin{lemma}[Normal-map criterion]\label{lem:core-normal}
If $N$ is admissible for $r$, then
\begin{equation}\label{eq:core-conclusion}
 \im s_{2r}(\Mat_{N+1}(K))=\Sl_{N+1}(K).
\end{equation}
Indeed the defining witness yields an attained nonzero rank-one
square-zero value at which the full differential is surjective.
At the particular border tuple used below, assuming $S$ invertible,
this differential condition is equivalent to injectivity of
$\Gam_{B,u}$.
\end{lemma}
\begin{proof}
Embed the fixed matrices and the final input as
\begin{equation}\label{eq:core-border}
 A_i=\begin{pmatrix}B_i&0\\0&0\end{pmatrix},\qquad
 Y=\begin{pmatrix}0&u\\0&0\end{pmatrix}.
\end{equation}
Only words ending in $Y$ can survive, because $YA_i=0$.
Since $k$ is odd, the value of the even alternator is
\begin{equation}\label{eq:core-border-value}
 s_{k+1}(A_1,\ldots,A_k,Y)
 =\begin{pmatrix}0&Su\\0&0\end{pmatrix}.
\end{equation}
If $u=0$, then $\Gam=L$ and its kernel contains all
$B_i$ (and also the identity), contradicting injectivity.
Hence $Su\ne0$. The value is nonzero, square zero and rank one.

To compute the differential, pair its values with a general
matrix using the nondegenerate trace form. Variations of $Y$
in the upper and lower borders give the full upper and lower
border directions by left and right multiplication by $S$ and
$-S$, respectively. Thus every normal is block diagonal.
The differential values have trace zero, so scalar matrices
are normals by \eqref{eq:even-trace}. Subtracting a scalar we
may write a normal uniquely as $\diag(P,0)$.

Variations of the old $N\times N$ block of $Y$ give
$\tr(P L(Q))$ for arbitrary $Q$. The map $L$ is
skew-adjoint for the trace pairing. To see this, collect the
cyclic classes in
$\tr s_{k+2}(B_1,\ldots,B_k,Q,P)$. The length
$k+2$ is odd, so a cyclic shift preserves sign; moving $Q$
past the $k$ frozen letters contributes a minus sign.
Consequently the trace equals $-(k+2)\tr(Q L(P))$.
Interchanging $P,Q$ in the alternating trace yields
\begin{equation}\label{eq:core-skew}
 \tr(P L(Q))=-\tr(Q L(P)).
\end{equation}
The old-block normal equations are therefore $L(P)=0$.

Now vary the $i$th $A_i$ in the lower border by a row
$c\in(K^N)^*$. For a nonzero word the unique upper-border
factor $Y$ must be immediately followed by this lower-border
variation. Remove $Y$ and rotate the remaining $k$ letters
cyclically. If $a$ of the original old letters followed the
adjacent pair, the removal has sign $(-1)^{a+1}$ and moving
the varied letter to its original position has sign $(-1)^a$.
The product is $-1$, and the normal pairing is exactly
\begin{equation}\label{eq:core-low-normal}
 -c\,C_i(P)u.
\end{equation}
Because all rows $c$ are available, these normal equations say
$C_i(P)u=0$ for each $i$. Therefore every normal modulo
scalars satisfies $\Gam_{B,u}(P)=0$. Conversely, in the pure
border tuple \eqref{eq:core-border}, these are all the normal
equations: other source-block variations yield only upper-border
outputs or zero. This also proves the asserted equivalence.
If $\Gam$ is injective, all normals are scalar, so the
full differential has image $\Sl_{N+1}$.
Lemma~\ref{lem:nilpotent-open} now proves
\eqref{eq:core-conclusion} as an \emph{attained} image equality.
\end{proof}

A complementary rectangular nonvanishing calculation records the
invertibility available for independently varying block-diagonal inputs.
\begin{lemma}[Mixed-rectangle nonvanishing]\label{lem:mixed-rect}
Let $a,b\ge1$ and $m<2(a+b)$. For an $m$-variable standard
polynomial, the free rectangular map
$\Mat_{a\times b}\to\Mat_{a\times b}$ obtained by freezing
$m-1$ block-diagonal inputs of sizes $a$ and $b$ is generically
invertible as the frozen blocks vary independently.
\end{lemma}
\begin{proof}
Choose integers $e,f\ge0$ with $e+f=m-1$,
$e<2a$ and $f<2b$; they exist since $m<2(a+b)$.
Specialize the first $e$ frozen labels to the $a$ block and
the remaining $f$ to the $b$ block. The only surviving words
put the left labels before the free rectangle and the right
labels after it, so the map is
\begin{equation}\label{eq:mixed-rect}
 X\longmapsto(-1)^f s_e(A)X s_f(B),\qquad s_0=I.
\end{equation}
Both factors admit invertible evaluations. Indeed, for degree
below the Amitsur--Levitzki threshold a standard polynomial
is a nonidentity. Expand a nonzero evaluation in the
clock-and-shift homogeneous basis of the corresponding matrix
algebra; by multilinearity one homogeneous tuple yields a
nonzero scalar multiple of an invertible basis element.
Thus the rectangular map is invertible at one specialization,
and its determinant is a nonzero polynomial.
\end{proof}

\begin{theorem}[Additivity of admissible core sizes]\label{thm:core-add}
For a fixed $r\ge1$, if $a,b\ge r$ are admissible core sizes,
then $a+b$ is admissible for $r$.
\end{theorem}
\begin{proof}
Choose witnesses $(B^{(a)},u_a)$ and $(B^{(b)},u_b)$
for the two sizes. For one parameter $t$ set
\begin{equation}\label{eq:core-add-tuple}
 B_i(t)=\diag(B_i^{(a)},tB_i^{(b)}),\qquad
 u=\binom{u_a}{u_b}.
\end{equation}
For every $t\ne0$, the diagonal endpoint blocks are
$S_a$ and $t^kS_b$ and hence are invertible. The old-block
$L$ and $C_i$ operators in the second block are respectively
scaled by $t^k$ and $t^{k-1}$, so their individual
admissibility kernels remain zero.

On the two off-diagonal rectangular spaces, the map $L_t$
in \eqref{eq:core-definitions} is polynomial in $t$. At $t=0$
the top-right rectangular map is $X\mapsto S_aX$,
because only words with the free variable last survive.
The bottom-left rectangular map is $X\mapsto -XS_a$,
because only words with the free variable first survive and
$k$ is odd. Both are invertible. Therefore their determinants
are nonzero polynomials in the \emph{same} scalar parameter.
As $K$ is infinite, select one common $t\in K^\times$
avoiding their zero sets.

At this fixed $t$, if $\Gam_{B(t),u}(P)=0$, the two mixed
$L_t$ equations force the off-diagonal blocks of $P$ to be
zero. The remaining $C_i(P)$ are block diagonal, and the
normal-map equations reduce to those for the two original
cores, with nonzero scalar factors on the second block.
Both diagonal blocks of $P$ vanish by the original
injectivity hypotheses. Thus $\Gam_{B(t),u}$ is injective,
and $\det s_k(B(t))\ne0$. This constructs a single
admissible witness of size $a+b$.
\end{proof}

\begin{corollary}[An infinite congruence family]\label{cor:congruence}
For every $r,q\ge1$,
\begin{equation}\label{eq:congruence}
 \im s_{2r}(\Mat_{qr+1}(K))=\Sl_{qr+1}(K).
\end{equation}
\end{corollary}
\begin{proof}
In Theorem~\ref{thm:even-seed}, set $n=r+1$, place the
exit variable last, and remove its sink coordinate. The old
$r\times r$ core has
\begin{equation}\label{eq:core-star}
 B=(E_{00},E_{01},E_{10},\ldots,E_{0,r-1},E_{r-1,0}),
 \qquad u=e_0,
\end{equation}
with the obvious singleton interpretation when $r=1$.
Counting closed trails through the root loop and the $r-1$
two-edge excursions gives
\begin{equation}\label{eq:core-star-endpoint}
 S=s_{2r-1}(B)
 =\diag(r!,(r-1)!,\ldots,(r-1)!),
\end{equation}
which is invertible in characteristic zero.
The full differential in Theorem~\ref{thm:even-seed} is onto
$\Sl_{r+1}$; the converse statement in
Lemma~\ref{lem:core-normal} shows that this old core is
admissible. For $r=1$ it follows also directly from
$S=1$ and $C_1(P)u=P$. Add $q$ copies of the admissible
size $r$ with Theorem~\ref{thm:core-add}, then apply
Lemma~\ref{lem:core-normal} in size $qr+1$.
\end{proof}

\section{Exact core certificates in degrees eight and ten}\label{sec:even-low}
For each $r$, the uniform construction supplies an admissible core
of size $r$. By additive transfer, to obtain all sizes $N\ge r$,
it suffices to establish admissibility for the finite range
\begin{equation}\label{eq:core-finite-reduction}
 r,r+1,\ldots,2r-1.
\end{equation}
Indeed every $N\ge r$ has the form $N=qr+s$,
$q\ge1$ and $0\le s<r$. If $s=0$, use $q$ copies of
$r$; otherwise use $q-1$ copies of $r$ and one copy of $r+s$.
Thus a finite collection of admissible cores extends to every
larger core order by Theorem~\ref{thm:core-add}.

For $r=4$ and $r=5$, the seven residual cases have exact
integer witnesses. For each case in the table, the corresponding
certificate contains an explicit ordered tuple $B$ of integer matrices,
an integer column $u$,
$\det S\ne0$, and the row indices and exact nonzero integer
determinant of an $N^2\times N^2$ submatrix of
$\Gam_{B,u}$. The last two columns give reductions of these full integer
minors modulo $101$ and $103$; their nonzero residues verify
nonvanishing in characteristic zero.
\begin{center}
\small
\begin{tabular}{@{}rrrcrr@{}}
\toprule
$r$ & $N$ & $2r$ & $\rk\Gam$ & minor $\bmod 101$ & minor $\bmod 103$\\
\midrule
4&5&8&25&48&57\\
4&6&8&36&36&25\\
4&7&8&49&75&58\\
5&6&10&36&52&34\\
5&7&10&49&64&45\\
5&8&10&64&45&60\\
5&9&10&81&85&7\\
\bottomrule
\end{tabular}
\end{center}
The certificate names are listed in
Appendix~\ref{app:even-verification}; they encode each pair $(r,N)$
from the table. For example,
\path{certificates/border_r4_N5.json} records $(r,N)=(4,5)$. The normal map has $N^2$
columns, so a nonzero $N^2$-minor proves its injectivity.
The invertible endpoint is separately recorded for each case.

\begin{theorem}[Degrees eight and ten without generic-nullity input]
\label{thm:even-eight-ten}
For every algebraically closed characteristic-zero field $K$,
\begin{equation}\label{eq:eight-ten}
\begin{aligned}
 \im s_8(\Mat_n(K))&=
 \begin{cases}\{0\},&n\le4,\\\Sl_n(K),&n\ge5,\end{cases}\\
 \im s_{10}(\Mat_n(K))&=
 \begin{cases}\{0\},&n\le5,\\\Sl_n(K),&n\ge6.\end{cases}
\end{aligned}
\end{equation}
\end{theorem}
\begin{proof}
\enlargethispage{2\baselineskip}
For $r=4$, the uniform seed gives admissible size $N=4$,
and the three finite residual sizes in
\eqref{eq:core-finite-reduction} are $N=5,6,7$.
For $r=5$, the uniform seed gives $N=5$, and the finite
residual sizes are $N=6,7,8,9$. Thus the seven exact
certificate witnesses together with Theorem~\ref{thm:core-add}
make every core size $N\ge r$ admissible, for $r=4,5$.
Lemma~\ref{lem:core-normal} yields the actual traceless image
for every order $n=N+1\ge r+1$. The other orders vanish by
Amitsur--Levitzki.
\end{proof}

Appendix~\ref{app:even-verification} describes how to check the
seven saved certificates. The extension to all matrix orders
uses Theorem~\ref{thm:core-add}.

\section{Generic nullity and the all-even classification}\label{sec:even-general}
We now deduce the remaining even degrees from two precise
external generic-kernel theorems. The argument extends the
approach to $s_6$ in Qiyue Tang's Appendix~A
\cite[Proposition~11]{Tang2026FiveSix}; in particular,
its use of a rank-one nilpotent, a generic central division
algebra and a projective eigenline incidence is directly indebted
to that source.

Write
\begin{equation}\label{eq:generic-op}
 \mathcal L_{A_1,\ldots,A_j}(P)
 =s_{j+1}(A_1,\ldots,A_j,P).
\end{equation}
The two inputs needed below are as follows. Brassil and
Reichstein \cite[Theorem~1.2]{BrassilReichstein2022} prove that
for $j$ even, $2\le j\le2N-2$, and generic $j$-tuples in
$\Mat_N$, the kernel of \eqref{eq:generic-op} has dimension
$j$ in characteristic zero. Their theorem is stated more
generally over an infinite field with specified characteristic
exclusions; only the characteristic-zero specialization is
used here. Lin, Chenhao Tang and Zhao
\cite[Theorem~1.2 and Remark~6.1]{LinTangZhao2026}
prove the complementary odd-frozen statement: for $N\ge3$,
odd $j$ with $3\le j\le2N-3$ and any characteristic-zero
field, the generic nullity is
\begin{equation}\label{eq:GN}
 \dim\ker\mathcal L_{A_1,\ldots,A_j}
 =\begin{cases}j+1,&N\text{ even},\\
 j+2,&N\text{ odd}.\end{cases}
\end{equation}
These generic-kernel dimensions provide the input for the
following attained-image argument.

\begin{theorem}[Even-standard classification from generic nullity]
\label{thm:even-complete}
For every algebraically closed field $K$ of characteristic zero
and all $r,n\ge1$,
\begin{equation}\label{eq:even-complete}
 \im s_{2r}(\Mat_n(K))=
 \begin{cases}\{0\},&n\le r,\\\Sl_n(K),&n\ge r+1.\end{cases}
\end{equation}
For $r\ge2$ and $n\ge r+2$, the proof uses the above external
Brassil--Reichstein and Lin--Tang--Zhao generic-nullity theorems.
\end{theorem}
\begin{proof}
If $n\le r$, then $2r\ge2n$, and the standard identity
forces the image to be zero. The case $r=1$ is the classical
commutator image theorem \cite{AlbertMuckenhoupt1957}.
For $r\ge2$ and $n=r+1$, the positive claim is
Theorem~\ref{thm:even-seed}.

Now let $r\ge2$ and $n\ge r+2$. Set
\begin{equation}\label{eq:generic-parameters}
 N=n-1,\qquad k=2r-1;
 \quad N\ge r+1\ge3,\quad 3\le k\le2N-3.
\end{equation}
For each $C_i$ of \eqref{eq:core-definitions}, the number
$k-1=2r-2$ of frozen inputs is even and satisfies
$2\le k-1\le2N-2$. Therefore both external nullity theorems
apply in precisely their stated ranges.

Take $k$ independent generic $N\times N$ matrices
$B_1,\ldots,B_k$ over $K$, and let $D$ denote the central
division algebra of generic matrices generated by them, with
center $F$ and degree $N$. Write $E$ for the field of their
independent matrix-entry indeterminates. The natural generic
matrix representation embeds $D$ in $\Mat_N(E)$ and splits it
over $E$; the center $F$ embeds in $E$ as scalar matrices.
The matrices $B_i$ are linearly independent over $F$,
because their vectorized columns are independent already over
$E$ (some $k\times k$ minor is a nonzero polynomial, as
$k\le N^2$). Set
\begin{equation}\label{eq:generic-S}
 S=s_k(B_1,\ldots,B_k)\in D.
\end{equation}
Since $k<2N$, the standard polynomial is not an identity
of $\Mat_N$; genericity makes $S\ne0$ and thus
$S\in D^\times$. In the generic matrix model, all of the
maps $L,C_i$ from \eqref{eq:core-definitions} are
$F$-linear operators on the $N^2$-dimensional space $D$.
Their ranks persist under splitting-field extension.

The Brassil--Reichstein theorem gives, simultaneously for
every $i$,
\begin{equation}\label{eq:generic-C-kernel}
 \ker C_i=\Span_F\{B_j:j\ne i\}.
\end{equation}
Indeed alternation gives the displayed $(k-1)$-dimensional
subspace, and the generic theorem gives exactly this nullity.
The generic open conditions for the finitely many $C_i$
intersect in the irreducible generic input space, so one set
of generic entries realizes all the ranks at once.
The odd-frozen theorem \eqref{eq:GN} gives
\begin{equation}\label{eq:generic-W}
 W=\ker L/\Span_F(B_1,\ldots,B_k),\qquad
 \dim_F W=\begin{cases}1,&N\text{ even},\\2,&N\text{ odd}.
 \end{cases}
\end{equation}
(The upper bound $\dim W\le2$ would suffice below.)

Alternation also gives
\begin{equation}\label{eq:C-on-B}
 C_i(B_j)=\delta_{ij}S.
\end{equation}
The $F$-linear maps
\begin{equation}\label{eq:generic-T}
 T_i:W\longrightarrow D/F1,
 \qquad[P]\longmapsto S^{-1}C_i(P)+F1
\end{equation}
are consequently well defined. They are \emph{jointly
injective}. In fact, if $C_i(P)=b_iS$ for every $i$, then
\eqref{eq:C-on-B} gives
$P-\sum_i b_iB_i\in\bigcap_i\ker C_i=\{0\}$;
the last equality follows from \eqref{eq:generic-C-kernel}
and the linear independence of all the $B_i$.
This joint injectivity persists after scalar extension.

Choose an $F$-linear section $W\to\ker L$ and let
$\bar F$ be an algebraic closure of $F$. Fix a splitting
identification $D\otimes_F\bar F\simeq\Mat_N(\bar F)$.
For a class $[P]\in W_{\bar F}$ the endomorphisms
$S^{-1}C_i(P)$ are defined by this section; changing the
section adds scalars to them, so their eigenlines are
independent of the section. We claim that there is a nonzero
column $u\in\bar F^N$ which is not a common eigenvector of
$S^{-1}C_i(P)$ for \emph{any} nonzero class $[P]$.

If $\dim W=0$ the claim is immediate. If $\dim W=1$, joint
injectivity means at least one of these matrices is nonscalar;
its eigenvectors form a finite union of proper eigenspaces,
which does not exhaust $\bar F^N$. The substantive case is
$\dim W=2$. In that case form the projective closed incidence
\begin{equation}\label{eq:generic-incidence}
 \mathcal I=\bigl\{([P],[u])\in
 \mathbb P(W_{\bar F})\times\mathbb P^{N-1}_{\bar F}:
 u\wedge S^{-1}C_i(P)u=0\quad(1\le i\le k)\bigr\}.
\end{equation}
The incidence records \emph{simultaneously} all kernel classes
and their common eigenlines.

At the generic point of the projective line
$\mathbb P(W_{\bar F})$, write $P=P_0+tP_1$ using an
$F$-basis of $W$. Joint injectivity of \eqref{eq:generic-T}
shows that some $S^{-1}C_i(P)$ is noncentral in
$D\otimes_FF(t)$. This algebra is a division algebra:
it embeds in the Laurent-series division ring $D((t))$,
with central indeterminate $t$, and as a finite-dimensional
domain over $F(t)$ every nonzero element is invertible.
A noncentral element in that division algebra has an
irreducible separable minimal polynomial of degree $a\ge2$
over $F(t)$; its degree divides $N$. Upon splitting, its
reduced characteristic polynomial is the $N/a$-th power of
that minimal polynomial. Its distinct eigenvalues all have
eigenspace dimension $N/a\le\lfloor N/2\rfloor$.
Consequently the generic fiber of \eqref{eq:generic-incidence}
has dimension at most $\lfloor N/2\rfloor-1$.

At every special geometric point $[P]$, the scalar-extended
joint injectivity still ensures that some
$S^{-1}C_i(P)$ is nonscalar. Its eigenspaces have dimension
at most $N-1$, and hence the corresponding fiber has
dimension at most $N-2$. Thus any component of
$\mathcal I$ dominating the projective line has dimension
at most $\lfloor N/2\rfloor$, and a component over a
point has dimension at most $N-2$. Since $N\ge3$,
\begin{equation}\label{eq:incidence-bound}
 \dim\mathcal I\le N-2<N-1=\dim\mathbb P^{N-1}.
\end{equation}
The projection to $\mathbb P^{N-1}$ is proper, so its image
is a proper closed subset. Choose $[u]$ outside that image.
This proves the claim for the \emph{same} $u$ and every
nonzero kernel class.

For this column, the normal map $\Gam_{B,u}$ is injective
after the chosen splitting. Indeed $\Gam(P)=0$ first implies
$P\in\ker L$. If its class in $W$ were nonzero, all
$C_i(P)u=0$ would make $[u]$ a common eigenline, contrary
to construction. Therefore
$P=\sum_i b_iB_i$; then
$C_i(P)u=b_iSu=0$ for all $i$ forces each $b_i=0$.
We have also retained $\det S\ne0$.

It remains to obtain a witness over the originally specified
field $K$, rather than merely over an abstract splitting of
$D$. Choose an algebraic closure $\Omega$ of the
original generic-entry field $E$ containing an $F$-embedding
of $\bar F$. The generic matrix representation and the
chosen split representation of
$D\otimes_F\Omega\simeq\Mat_N(\Omega)$ differ by an inner
automorphism, by Skolem--Noether. Transport the column $u$
under that conjugacy. Equivariance of $S,L,C_i$ shows that
$\det S\ne0$ and $\rk\Gam=N^2$ hold in the original
generic-entry representation over $\Omega$.

\enlargethispage{1\baselineskip}
Both assertions mean that the product of $\det S$ with a
suitable $N^2\times N^2$ minor of $\Gam$ is a \emph{nonzero
polynomial} in the generic entries of $B$ and the coordinates
of $u$, with coefficients in $K$. Because $K$ is infinite,
there is one specialization of all these coordinates to $K$
at which this polynomial is nonzero. This gives an admissible
core of size $N=n-1$ over $K$ itself, in the sense of
\eqref{eq:core-admissible}. By Lemma~\ref{lem:core-normal},
the full polynomial map $s_{2r}$ attains a nonzero rank-one
nilpotent and therefore has actual image $\Sl_n(K)$.
This proves \eqref{eq:even-complete} in the remaining range.
\end{proof}



\begin{proof}[Proof of Theorem~\ref{thm:main}]
The odd assertion is proved in Section~\ref{sec:assembly}. The even
assertion follows from Theorem~\ref{thm:even-complete}, including the
identity and low-order cases. Both arguments establish equality of
sets of individual values: by explicit slices in odd degree and by
the attained-nilpotent criterion in even degree.
\end{proof}

\section*{Use of AI}
OpenAI's GPT-6 Pro assisted substantially in the mathematical
investigation, development of proofs, exact computation, source
assessment and manuscript preparation. The author is responsible for
the mathematical claims, attribution and final presentation.

\appendix
\section{Exact checks for odd standard polynomials}\label{app:verification}
The programs \path{code/verify_standard.py} and
\path{code/graded.py} provide finite exact checks of the odd
path, cycle, periodic and exceptional-order calculations.
The first uses SymPy and evaluates alternating products by a
full-matrix subset recurrence and by enumerating nonzero
matrix-unit Euler trails. The second computes cyclic-homogeneous
slice restrictions by a graded subset recurrence.

For fixed inputs $A_1,\ldots,A_m$, let $D(S)$ denote the alternating
sum over the variables indexed by $S\subseteq\{1,\ldots,m\}$.
The full-matrix recurrence is
\begin{equation}\label{eq:dp}
 D(\varnothing)=I,\qquad
 D(S)=\sum_{j\in S}(-1)^{|S|-\operatorname{pos}_S(j)}
 D(S\setminus\{j\})A_j,
\end{equation}
where $\operatorname{pos}_S(j)$ is the position of $j$ in the
increasing listing of $S$. For homogeneous cyclic inputs
$A_j=\diag(w_j)P^{g_j}$, the graded calculation tracks the
offset of each product and constructs the restriction matrix
directly. The two calculations compare every column
$L(E_{i,i+h\bmod n})$ for the specified exceptional cases;
the generated JSON also includes integer determinants,
finite trail counts, periodic pairs and metadata for sample Jordan
constructions. The examples cover orders $3$ through $8$, with
direct full-alternator evaluation checked for orders $3$ through $5$.

With Python 3.10 or later and the packages in
\path{requirements.txt}, run from the manuscript directory:
\begin{samepage}
\begin{verbatim}
python -m pip install -r requirements.txt
python code/verify_standard.py --output standard_rerun.json
python code/verify_standard.py --extended \
  --output standard_extended_rerun.json
\end{verbatim}
\end{samepage}
Compare the parsed outputs with
\path{certificates/standard_certificates.json} and
\path{certificates/standard_certificates_extended.json}, respectively.
The tests use assertions, so \texttt{python -O} must not be used.
The corresponding finite ranges are described in
\path{REPRODUCIBILITY.md}.

\section{Even-standard integer certificates}\label{app:even-verification}
The seven even-core certificates, indexed by their pairs $(r,N)$,
are the following files:
\begin{center}
\small
\begin{tabular}{@{}l@{}}
\path{certificates/border_r4_N5.json}\\
\path{certificates/border_r4_N6.json}\\
\path{certificates/border_r4_N7.json}\\
\path{certificates/border_r5_N6.json}\\
\path{certificates/border_r5_N7.json}\\
\path{certificates/border_r5_N8.json}\\
\path{certificates/border_r5_N9.json}
\end{tabular}
\end{center}
Each file contains integer matrices $B_1,\ldots,B_{2r-1}$,
a column $u$, the nonzero integer determinant of
$S=s_{2r-1}(B)$, and a specified $N^2\times N^2$ minor of
the normal map $\Gam_{B,u}$ with its integer determinant.
The nonzero determinant of $S$ and minor of $\Gam$ establish
admissibility \eqref{eq:core-admissible} for that core size.

The constructor \path{code/verify_border_even.py} uses an
alternating-product subset recurrence and replaced-input
columns, with direct word enumeration for smaller cases.
The saved-input checker \path{code/verify_saved_border.py}
uses a different first-letter recurrence to rebuild every
normal-map column; it checks $\det S$, the bordered rank-one
square-zero value and the recorded minor modulo $101$ and
$103$. A nonzero residue of the integer minor establishes
its nonvanishing over characteristic-zero fields.
The seed checker \path{code/verify_minimal_even.py} tests the
signed matrix-unit formulas and full differential in specified
finite sizes.

For example, run:
\begin{verbatim}
python code/verify_saved_border.py \
  --input certificates/border_r4_N5.json \
  --output checked_border_r4_N5.json
python code/verify_minimal_even.py \
  --max-n 8 --rank-max-n 6 \
  --output minimal_even_rerun.json
\end{verbatim}
The first command applies likewise to the other six certificate
files; \path{REPRODUCIBILITY.md} gives the full list and
constructor command. These integer checks establish the listed
finite admissible cores. Their extension to infinitely many
matrix orders is the symbolic additive transfer
of Theorem~\ref{thm:core-add}, while the general even-degree
argument uses the cited generic-nullity theorems.

\bibliographystyle{amsplain}
\urlstyle{same}
{\raggedright\interlinepenalty=10000
\bibliography{references}
}
\end{document}
